\subsection{矩阵计算。}

本小节中，我们打算推广第 II 章，第 6 节中引入的矩阵演算，并用它来转述本节中已证明的某些结果。

I. --- 设 I 与 K 是两个有限指标集，H 是一个非空集合，\(M = (m_{ik})_{(i,k) \in \mathbb{I} \times \mathbb{K}}\) 是 H 上的一个矩阵（第 II 章，第 6 节，第 1 小节，定义 1）。

称为 \(M\) 的转置（记作 \({}^t M\) ）的是矩阵 \((m_{ki}')_{(k,i) \in \mathbb{K} \times \mathbf{I}}\) ，它满足 \(m_{ki}' = m_{ik} ((i, k) \in \mathbf{I} \times \mathbf{K})\) 。显然有

\[
^ {t} (^ {t} M) = M.\tag{39}
\]

这推广了第 II 章，第 6 节，第 6 小节中引入的概念。

设 H 是一个交换群（记法为加法）。以 I 与 K 为指标集的 H 上矩阵的集合具有交换群结构，因为它是 I × K 到 H 的映射的集合。这个群用加法记之。

设 H′、H″ 是两个非空集合，H 是一个交换群（记法为加法），\(f:(h',h'')\to h'h''\) 是 \(\mathbf{H}'\times\mathbf{H}''\) 到 H 中的一个映射。给定两个矩阵

\[
M ^ {\prime} = (m _ {i k} ^ {\prime}) _ {(i, k) \in \mathrm{I} \times \mathrm{K}}, \quad M ^ {\prime \prime} = (m _ {k l} ^ {\prime \prime}) _ {(k, l) \in \mathrm{K} \times \mathrm{L}}
\]

（分别在 H′ 与 H″ 上），使得 M′ 的列指标集 K 等于 M″ 的行指标集；称 M′ 与 M″ （依 f ）的积，并记作 M′M″ ，为矩阵

\[
M ^ {\prime}. M ^ {\prime \prime} = (\sum_ {k \in \mathbf {K}} m _ {i k} ^ {\prime} m _ {k l} ^ {\prime \prime}) _ {(i, l) \in \mathbf {I} \times \mathbf {L}}\tag{40}
\]

（在 H 上）。这推广了第 II 章，第 6 节，第 4 小节中引入的概念。若 \(\mathbf{H}' = \mathbf{H}'' = \mathbf{H}\) 且 H 是一个环，则除非另有明确说明，积 \(M'M''\) 将“在 H 中”计算，也就是按照映射 \((x, y) \to xy\) 计算。当 \(\mathbf{H}'\) 与 \(\mathbf{H}''\) 是交换群（记法为加法）且 \(f\) 是双线性映射时，有

\[
\left\{ \begin{array}{l} (M ^ {\prime} + M _ {1} ^ {\prime}) M ^ {\prime \prime} = M ^ {\prime} M ^ {\prime \prime} + M _ {1} ^ {\prime} M ^ {\prime \prime}, \\ M ^ {\prime} (M ^ {\prime \prime} + M _ {1} ^ {\prime \prime}) = M ^ {\prime} M ^ {\prime \prime} + M ^ {\prime} M _ {1} ^ {\prime \prime}, \end{array} \right.\tag{41}
\]

其中 \(M'\) 、\(M_1'\) 是 \(H'\) 上的矩阵，\(M''\) 、\(M_1''\) 是 \(H''\) 上的矩阵，且写出的和与积都假定有定义。设 \(M'\) 、\(M''\) 是集合 \(H'\) 、\(H''\) 上的矩阵，\(f^0\) 是 \(H'' \times H'\) 到 H 中的、由 \((h'', h') \to h' h''\) 定义的映射；则有

\[
{ } ^ { t } ( M ^ { \prime } M ^ { \prime \prime } ) = { } ^ { t } M ^ { \prime \prime } . { } ^ { t } M ^ { \prime }\tag{42}
\]

其中第一个（相应地第二个）成员中的积按照 \(f\) （相应地 \(f^{0}\) ）计算。

在 \(H' = H'' = H\) 是一个环的情形，我们重新得到第 II 章，第 6 节，第 6 小节的公式 (12)。

设 A 是一个环，J 是 A 的一个反自同构。对 A 上每个矩阵 \(M = (m_{ik})\) ，我们将用 \(M^{J}\) 表示矩阵 \((m_{ik}^{J})\) 。设 \(M_{1}\) 、\(M_{2}\) 是 A 上两个矩阵，使得 \(M_{1}M_{2}\) 有定义。由于 J 是 A 到反环 \(A^{0}\) 上的一个同构，有 \((M_{1}M_{2})^{J} = M_{1}^{J}\) . \(M_{2}^{J}\) ，其中第一个（相应地第二个）成员在 A （相应地 \(A^{0}\) ）中计算。鉴于 (42) 与 (39)，这给出

\[
(M _ {1} M _ {2}) ^ {\mathrm{J}} = ^ {t} (^ {t} M _ {2} ^ {\mathrm{J}}. ^ {t} M _ {1} ^ {\mathrm{J}})\tag{43}
\]

其中两个成员都在 A 中计算。

设 \(H_{1}\) 、\(H_{2}\) 、\(H_{3}\) 、\(H_{12}\) 、\(H_{23}\) 与 H 是交换群（记法为加法），\(f_{12}: H_{1} \times H_{2} \to H_{12}\) 、\(f_{23}: H_{2} \times H_{3} \to H_{23}\) 、\(f_{3}: H_{12} \times H_{3} \to H\) 、\(f_{1}: H_{1} \times H_{23} \to H\) 是映射，并设 \(M_{1}\) 、\(M_{2}\) 、\(M_{3}\) 分别是 \(H_{1}\) 、\(H_{2}\) 、\(H_{3}\) 上的矩阵。若有 \(f_{3}(f_{12}(x_{1}, x_{2}), x_{3}) = f_{1}(x_{1}, f_{23}(x_{2}, x_{3}))\) （对一切 \(x_{i} \in H_{i}\) ，i = 1, 2, 3），则积 \((M_{1}M_{2})M_{3}\) 与 \(M_{1}(M_{2}M_{3})\) （分别按照 \(f_{12}\) 、\(f_{3}\) 、\(f_{23}\) 与 \(f_{1}\) 计算）若有定义便相等；将把它们记作 \(M_{1}M_{2}M_{3}\) 。当 \(H_{1} = H_{2} = H_{3} = H_{12} = H_{23} = H\) 、H 是一个环且 \(f_{12}\) 、\(f_{23}\) 、\(f_{3}\) 、\(f_{1}\) 都等于映射 \((x, y) \to xy\) 时，前一条件表示后者的结合律，因此是满足的。对多于三个因子的积，将作类似的约定。

设 A、B 是两个环，\(M = (m_{ik})_{(i,k) \in I \times K}\) 与 \(M' = (m_{ik}')_{(i,k) \in I \times K}\) 是 (A, B)-双模 G 上的两个矩阵（第 1 小节）。若对每个元素在 A 中的单行矩阵 \(L = (a_i)_{i \in I}\) 与每个元素在 B 中的单列矩阵 \(C = (b_k)_{k \in K}\) ，都有 \(L.M.C = L.M'.C\) （积按照定义 G 的 (A, B)-双模结构的映射计算），则矩阵 \(M\) 与 \(M'\) 相等。事实上，若取 \(a_i = 1\) ， \(a_s = 0\) （ \(s \neq i\) ）， \(b_k = 1\) ， \(b_t = 0\) （ \(t \neq k\) ），则纯量矩阵 \(L.M.C\) 与 \(L.M'.C\) 分别等于 \(m_{ik}\) 与 \(m_{ik}'\) 。

\begin{enumerate}
\def\labelenumi{\Roman{enumi}.}
\setcounter{enumi}{1}
\tightlist
\item
  考虑一个环 A 与一个（右或左）A-模 E，它具有有限基 \((e_{i})_{i\in I}\) 。对 E 的每个元素 x，由分量 \(x_{i}\) （ \(i\in I\) ，即 x 关于 \((e_{i})\) 的分量）构成的单列矩阵称为 x 关于基 \((e_{i})\) 的矩阵，记作 \(M(x)\) 或 x（参看第 II 章，第 6 节，第 4 小节）；在计算中，为记起指标 i 是行指标，给它添上一个只能取一个值的列指标，并用 \((x_{i0})\) 来记矩阵 \(M(x)\) ，将是方便的。
\end{enumerate}

现在考虑两个（左或右）A-模 E 与 F，分别具有有限基 \((e_{i})_{i\in I}\) 与 \((f_{k})_{k\in K}\) ；设 \((f_{k}^{*})\) 是 \(F^{*}\) 的与 \((f_{k})\) 对偶的基。我们将在下列四种情形中，定义 E 到 F 中的映射 u 关于这些基的矩阵：

(D) E 与 F 是右模，且 u 是 A-线性的；

(G) E 与 F 是左模，且 u 是 A-线性的；

(GD) E 是左模，F 是右模，A 配备了一个反自同构 J，u 是 Z-线性的且满足 \(u(ax) = u(x)a^{J}\) （ \(a \in A, x \in E\) ）（换句话说，u 是 \(E^{J}\) 到 F 中的 A-线性映射（第 2 小节，定义 5））。

(DG) E 是右 A-模，F 是左 A-模，A 配备了一个反自同构 J，u 是 Z-线性的且满足 \(u(xa) = a^{J}u(x)\) （ \(x \in E, a \in A\) ）（换句话说，u 是 \(E^{J}\) 到 F 中的一个 A-线性映射）。

在这四种情形的每一种中，映射 u 的矩阵按定义是使得下式成立的矩阵 \((u_{ki})_{(k,i)\in\mathbb{K}\times\mathbb{I}}\) ：

\[
u _ {k i} = \left\langle u (e _ {i}), f _ {k} ^ {*} \right\rangle .\tag{44}
\]

在情形 (D)，这个定义与第 II 章，第 6 节，第 3 小节给出的定义一致。在这些条件下，矩阵 \(M(u(x))\) （E 的元素 \(x\) 的像的矩阵）由下列公式给出：

\[
M (u (x)) = M (u). M (x)\tag{45 D}
\]

\[
{ } ^ { t } M ( u ( x ) ) = { } ^ { t } M ( x ) . { } ^ { t } M ( u )\tag{45 G}
\]

\[
M (u (x)) = M (u). M (x) ^ {\mathrm{J}}\tag{45 GD}
\]

\[
{ } ^ { t } M ( u ( x ) ) = { } ^ { t } M ( x ) ^ { j } . { } ^ { t } M ( u ) .\tag{45 DG}
\]

例如验证 (45 DG)，其余验证类似且略更容易。设 \(x = \sum_{i} e_i x_{io}, u(x) = \sum_{k} y_{ko} f_k\) ；我们有 \(u(x) = u(\sum_{i} e_i x_{io}) = \sum_{i} x_{io}^{\mathbf{J}} u(e_i) = \sum_{i,k} x_{io}^{\mathbf{J}} u_{ki} f_k\) ；由此得 \(y_{ko} = \sum_{i} x_{io}^{\mathbf{J}} u_{ki}\) ；为使两个指标 \(i\) 彼此靠近，考虑转置矩阵 \(^t M(x) = (x_{oi}')\) （其中 \(x_{oi}' = x_{io}\) ）与 \(^t M(u) = (u_{ik}')\) （其中 \(u_{ik}' = u_{ki}\) ）；于是有 \(y_{ko} = \sum_{i} x_{oi}' u_{ik}'\) ；由于右端是以 \(k\) 为指标的元素（属于单行矩阵 \(^t M(x)^{\mathbf{J}}\) . \(^t M(u)\) ），公式 (45 DG) 得到验证。

注记. -- 1) 当 A 交换时，借助公式 \(t(M'M'') = 'M'''.'M'\) （参看 (42)），(45 G) 化为 (45 D)，(45 DG) 化为 (45 GD)，其中两端的成员这里都在 A 中计算。2) 设 E、F、G 是三个具有有限基的左模，\(u: \mathrm{E} \to \mathrm{F}\) 、\(\nu: \mathrm{F} \to \mathrm{G}\) 是 A-线性映射。由 (45 G) 得出

\[
{ } ^ { t } M ( \nu \circ u ) = { } ^ { t } M ( u ) . { } ^ { t } M ( \nu ) .\tag{46}
\]

事实上，对任意 \(x \in E\) ，我们有

\[
\begin{array}{r l} ^ {t} M (x). ^ {t} M (\nu \circ u) & = ^ {t} M (\nu (u (x))) = ^ {t} M (u (x)). ^ {t} M (\nu) \\ & = ^ {t} M (x). ^ {t} M (u). ^ {t} M (\nu), \end{array}
\]

由此得 (46)。

回顾在右模的情形，我们有

\[
M (\nu \circ u) = M (\nu) M (u).
\]

\begin{enumerate}
\def\labelenumi{\Roman{enumi}.}
\setcounter{enumi}{2}
\tightlist
\item
  从现在起，我们用 A 表示一个环，用 B 表示一个环（相应地一个配备了反自同构 J 的环，对此我们令 \(J' = J^{-1}\) ），用 E 表示具有有限基 \((e_i)_{i \in I}\) 的一个左 A-模，用 F 表示具有有限基 \((f_k)_{k \in K}\) 的一个右（相应地左）B-模。我们用 \((e_i^*)\) 与 \((f_k^*)\) 表示 \(E^*\) 与 \(F^*\) 的对偶基。除非另有明确说明，所考虑的矩阵都关于这些基取定。
\end{enumerate}

设 G 是一个 (A, B)-双模（第 1 小节），Φ 是 E × F 到 G 中的一个双线性（相应地关于 J 的右半双线性）映射，\(R = (\Phi(e_i, f_k))\) 是 Φ 的矩阵。于是，对 \(x \in E\) 与 \(y \in F\) ，第 1 小节的公式 (6)（相应地第 2 小节的 (8)），鉴于上述约定，写成

\[
\Phi (x, y) = ^ {t} M (x). R. M (y) \quad (\text { resp. } \Phi (x, y) = ^ {t} M (x). R. M (y) ^ {\mathrm{J}}),\tag{47}
\]

其中积按照定义 (A, B)-双模 G 的结构的映射计算；特别地，若 A = B = G（这时 Φ 是一个型），则积在 A 中计算。

设 E′ 是具有有限基 \((e_{s}^{\prime})_{s\in S}\) 的一个左 A-模，F′ 是具有有限基 \((f_{t}^{\prime})_{t\in T}\) 的一个右（相应地左）A-模，\(u:E\to E^{\prime}\) 与 \(\nu:F\to F^{\prime}\) 是 A-线性映射，\(\Phi^{\prime}\) 是 \(E^{\prime}\times F^{\prime}\) 到 G 中的一个双线性（相应地关于 J 的右半双线性）映射。用 \(\Phi\) 表示 \(\Phi^{\prime}\) （关于 u 与 \(\nu\) ）的原像，用 U、V、R、\(R^{\prime}\) 表示 u、\(\nu\) 、\(\Phi\) 、\(\Phi^{\prime}\) 关于所考虑的基的矩阵。于是有

\[
R = ^ {t} U. R ^ {\prime}. V \quad (\text { resp. } R = ^ {t} U. R ^ {\prime}. V ^ {J}),\tag{48}
\]

其中积如 (47) 中那样计算。事实上，无论 \(x \in \mathbf{E}\) 与 \(y \in \mathbf{F}\) 如何，按定义有 \(\Phi(x, y) = \Phi'(u(x), v(y))\) ，于是由 (47)，

\[
\begin{array}{c} ^ {t} M (x). R. M (y) = ^ {t} M (u (x)). R ^ {\prime}. M (\nu (y)) \\ (\text {resp.} ^ {t} M (x). R. M (y) ^ {\mathrm{J}} = ^ {t} M (u (x)). R ^ {\prime}. M (\nu (y)) ^ {\mathrm{J}}); \end{array}
\]

由 (45 G) 与 (45 D)（相应地 (45 G)）以及 (43)，导出

\[
\begin{array}{r l} ^ {t} M (x). R. M (y) & = ^ {t} M (x). ^ {t} U. R ^ {\prime}. V. M (y) \\ (\text {resp.} ^ {t} M (x). R. M (y) ^ {J} & = ^ {t} M (x). ^ {t} U. R ^ {\prime}. ^ {t} (^ {t} M (y). ^ {t} V) ^ {J} \\ & = ^ {t} M (x). ^ {t} U. R ^ {\prime}. V ^ {J}. M (y) ^ {J}); \end{array}
\]

这就证明了我们的断言。

\begin{enumerate}
\def\labelenumi{\Roman{enumi}.}
\setcounter{enumi}{3}
\tightlist
\item
  --- 这里假设环 A 与 B 相等，并用 \(\Phi\) 表示 E × F 上的一个双线性（相应地关于 J 的右半双线性）型，用 \(R\) 表示其矩阵。我们来计算映射 \(s_{\Phi}\) 与 \(d_{\Phi}\) （与 \(\Phi\) 相伴者）的矩阵，为简单起见将把它们记作 \(s\) 与 \(d\) 。由于由第 6 小节的 (23)（相应地由第 6 小节的 (24)）有 \(\Phi(x, y) = \langle y, s(x) \rangle = \langle x, d(y) \rangle\) （相应地 \(\Phi(x, y) = \langle x, d(y) \rangle = \langle y, s(x) \rangle^{J}\) ），便有 \(\Phi(e_i, f_k) = \langle f_k, s(e_i) \rangle = \langle e_i, d(f_k) \rangle\) （相应地 \(\Phi(e_i, f_k) = \langle e_i, d(f_k) \rangle = \langle f_k, s(e_i) \rangle^{J}\) ），于是，由 (44) 并由于 ( \(e_i\) ) 是 ( \(e_i^*\) ) 的对偶基、( \(f_k\) ) 是 ( \(f_k^*\) ) 的对偶基：
\end{enumerate}

\[
M (d) = R,   M (s) = ^ {t} R \quad (\text { resp. }   M (d) = R,   M (s) = ^ {t} R ^ {J ^ {\prime}}).\tag{49}
\]

注记. --- 1) 当 A 是一个体时，线性映射 s 与 d 有相同的秩。我们在这里看到，它们的矩阵 \(M(s)\) 与 \(M(d)\) 有相同的秩；事实上，A 上的矩阵与其转置有相同的秩（第 II 章，第 6 节，第 7 小节，命题 3），且当 \(\Phi\) 是半双线性时，\(R\) 在 A 上的秩与 \(^tR\) 在 A \(^0\) 上的秩相等（同处）以及 J′ 是 A \(^0\) 到 A 上的同构这一事实，蕴涵 \(R\) 与 \(^tR^{J'}\) 在 A 上的秩相等。

\begin{enumerate}
\def\labelenumi{\arabic{enumi})}
\setcounter{enumi}{1}
\tightlist
\item
  若 M 与 N 是具有有限基 \((m_{i})\) 与 \((n_{k})\) 的右 A-模，\(\Phi\) 是 \(M \times N\) 上的一个（关于 J 的）左半双线性型（第 6 小节，注记），s 与 d 是它的相伴映射，\(R = (\Phi(m_{i}, n_{k}))\) 是它的矩阵，则第 6 小节的公式 (26) 表明有
\end{enumerate}

\[
M (d) = R ^ {J ^ {\prime}}, \quad M (s) = ^ {t} R.
\]

现在设与 \(\Phi\) 相伴的映射 s 与 d 是双射，而来计算矩阵 \(\hat{R}\) （\(\Phi\) 的逆型的矩阵，第 7 小节）。当 \(\Phi\) 是双线性型时，\(\Phi\) 是 \(\hat{\Phi}\) 关于线性映射 \(s: \mathrm{E} \to \mathrm{F}^*\) 与 \(d: \mathrm{F} \to \mathrm{E}^*\) 的原像；于是，鉴于 (48) 与 (49)，有 \(R = R\) . \(\hat{R}\) . \(R\) ，从而，由于 \(R\) 可逆（ \(d\) 是双射），\(\hat{R} = R^{-1}\) 。这个公式可推广到 \(\Phi\) 是半双线性型的情形，因为，若把 \(\Phi\) 看作 \(\mathrm{E} \times \mathrm{F}^{\mathrm{j}}\) 上的双线性型，并把 \((\mathrm{F}^{\mathrm{j}})^{*}\) 与 \((\mathrm{F}^{*})^{\mathrm{j}}\) 等同起来（参看第 9 小节），则这个双线性型的逆型与 \(\hat{\Phi}\) （看作 \((\mathrm{F}^{*})^{\mathrm{j}} \times \mathrm{E}^{*}\) 上的双线性型）一致。在两种情形下，\(\Phi\) 的逆型的矩阵都是 \(\Phi\) 的矩阵的逆。

最后，设 E′ 是一个左 A-模，F′ 是一个右（相应地左）A-模，两者都具有有限基 \((e_{s}^{\prime})\) 与 \((f_{t}^{\prime})\) ；设 \(\Phi^{\prime}\) 是 \(E^{\prime} \times F^{\prime}\) 上的一个双线性型（相应地关于 J 的半双线性型），\(R^{\prime}\) 是它的矩阵。设 \(s_{\Phi}\) 与 \(d_{\Phi}\) 是双射。设 \(u : E \to E^{\prime}\) 与 \(\nu : F \to F^{\prime}\) 是线性映射，\(u^{*} : F^{\prime} \to F\) 与 \(\upsilon^{*} : E^{\prime} \to E\) 是它们的伴随（第 8 小节，命题 7）；用 U、V、\(U^{*}\) 、\(V^{*}\) 表示 u、\(\upsilon\) 、\(u^{*}\) 、\(\upsilon^{*}\) 关于给定基的矩阵。于是有

\[
U ^ {*} = R ^ {- 1}. ^ {t} U. R ^ {\prime}, \quad {} ^ {t} V ^ {*} = R ^ {\prime}. V. R ^ {- 1}\tag{50}
\]

\[
(\text { resp. } U ^ {* \mathrm{J}} = R ^ {- 1}. ^ {t} U. R ^ {\prime}, ^ {t} V ^ {*} = R ^ {\prime}. V ^ {\mathrm{J}}. R ^ {- 1}).
\]

事实上，无论 \(x \in E\) 与 \(y \in F'\) 如何，有 \(\Phi'(u(x), y) = \Phi(x, u^*(y))\) （第 8 小节，定义 10）。于是，当 \(\Phi\) 是双线性型时，鉴于 (47)，有 \(^t M(u(x)). R'. M(y) = ^t M(x). R. M(u^*(y))\) ；这给出，鉴于 (45 G) 与 (45 D)，\(^t M(x). ^t U. R'. M(y) = ^t M(x). R. U^*. M(y)\) ，从而 \(^t U. R' = R. U^*\) 以及第一个所宣布的公式，因为 \(d\) 是双射，\(R\) 可逆。当 \(\Phi\) 是半双线性型时，(47) 与 (45 G) 给出 \(^t M(x). ^t U. R'. M(y)^J = ^t M(x). R. (^{t}M(y). ^{t}U^*)^{J}\) ；而由 (43)，我们有 \((^{t}M(y). ^{t}U^*)^{J} = ^{t}(^{u}U^{*J}. ^{u}M(y)^{J})\) ，从而 \((^{t}M(y). ^{t}U^*)^{J} = U^{*J}. M(y)^{J}\) ；因此导出 \(^t M(x). ^{t}U. R'. M(y)^{J} = ^{t}M(x). R. U^{*J}. M(y)^{J}\) ，从而 \(^t U. R' = R. U^{*J}\) ，以及 \(U^{*J} = R^{-1}. ^{t}U. R'\) 。关于 \(V^*\) 的公式的验证是类似的。

习题. --- 1) 设 A 是一个交换域，E 是 A 上的一个向量空间，具有可数无穷基 \((e_n)_{n \geqslant 1}\) 。如下定义 E 上的一个双线性型 \(\Phi\) ：令 \(\Phi(e_{i+1}, e_i) = 1\) （对 \(i \geqslant 1\) ），\(\Phi(e_k, e_j) = 0\) （对 \(k \neq j + 1\) 与 \(j \geqslant 1\) ）。证明线性映射 \(d_\Phi\) （与 \(\Phi\) 右相伴者）是单射，但线性映射 \(s_\Phi\) （与 \(\Phi\) 左相伴者）不是单射。

\begin{enumerate}
\def\labelenumi{\arabic{enumi})}
\setcounter{enumi}{1}
\item
  设 E 是 Z 与 Z/(2) 的 Z-模直和，E* 是它的对偶（同构于 Z）。证明双线性型 \((x, x') \to \langle x, x' \rangle\) （在 \(E \times E^*\) 上）使得右相伴线性映射是单射，而左相伴线性映射不是。
\item
  给出一个双线性型 \(\Phi\) 的例子，它定义在两个向量空间的积 \(\mathbf{E} \times \mathbf{F}\) 上，使得 \(d_{\Phi}\) 是双射，而 \(s_{\Phi}\) 是单射但非双射（取 \(\mathbf{E}\) 为无限维，且 F 等于 \(\mathbf{E}^*\) ，即 \(\mathbf{E}\) 的对偶；参看第 II 章，第 5 节，习题 3）。
\item
  设 A 是一个配备了反自同构 J 的环，E 是一个左 A-模，G 是一个 (A, A)-双模，\(\Phi\) 是 \(\mathrm{E} \times \mathrm{E}\) 到 G 中的一个（关于 J 的右）半双线性映射。证明恒等式（其中 \(\mathrm{Q}(x) = \Phi(x, x)\) ）：
\end{enumerate}

\[
\begin{array}{r l} & {2 \Phi (x, y) (\mu^ {\mathrm{J}} \lambda^ {\mathrm{J}} - \lambda^ {\mathrm{J}} \mu^ {\mathrm{J}}) = \mathrm{Q} (x - \mu \lambda y) - \mathrm{Q} (x + \mu \lambda y) + \mu \mathrm{Q} (x + \lambda y)} \\ & {- \mu \mathrm{Q} (x - \lambda y) + \mathrm{Q} (x + \mu y) \lambda^ {\mathrm{J}} - \mathrm{Q} (x - \mu y) \lambda^ {\mathrm{J}} + \mu \mathrm{Q} (x - y) \lambda^ {\mathrm{J}} - \mu \mathrm{Q} (x + y) \lambda^ {\mathrm{J}}.} \end{array}
\]

\begin{enumerate}
\def\labelenumi{\arabic{enumi})}
\setcounter{enumi}{4}
\tightlist
\item
  设 K 是一个特征为 2 的交换域，A 是 K 的一个可分二次扩张；我们有 \(\mathrm{A} = \mathrm{K}(\theta)\) ，其中 \(\theta\) 是 \(\mathrm{X}^2 +\mathrm{X} + \beta\) （K{[}X{]} 的一个不可约多项式）的根，而 A 的异于恒同的 K-自同构 J 满足 \(\theta^{j} = \theta +1\) （第 V 章，第 11 节，习题 8）。证明：若 E 与 G 是 A 上的向量空间，\(\Phi\) 是 \(\mathrm{E}\times \mathrm{E}\) 到 G 中的一个（关于 J 的）半双线性映射，则令 \(\mathrm{Q}(x) = \Phi (x,x)\) 后，有
\end{enumerate}

\[
\Phi (x, y) = \mathrm{Q} (\theta x + y) - \beta \mathrm{Q} (x) - \mathrm{Q} (y) - (\theta + 1) (\mathrm{Q} (x + y) - \mathrm{Q} (x) - \mathrm{Q} (y)).
\]

\begin{enumerate}
\def\labelenumi{\arabic{enumi})}
\setcounter{enumi}{5}
\tightlist
\item
  设 A 是一个体，E 是 A 上的一个向量空间，\(\Phi\) 是 E 上的一个半双线性型，\(u\) 是 E 的一个自同态。
\end{enumerate}

\begin{enumerate}
\def\labelenumi{\alph{enumi})}
\item
  存在且仅存在 E 的一个自同态 \(u^*\) ，使得 \(\Phi(u(x), y) = \Phi(x, u^*(y))\) （对 E 中的 \(x, y\) ），必须且只需 \(d_\Phi\) 是单射且 \(^t u(d_\Phi(E)) \subset d_\Phi(E)\) 。
\item
  给出一个例子，其中 E 是无限维且 \(d_{\Phi}\) 是单射，但 \(u(d_{\Phi}(E))\) 不包含于 \(d_{\Phi}(E)\) 。
\end{enumerate}

\begin{enumerate}
\def\labelenumi{\arabic{enumi})}
\setcounter{enumi}{6}
\tightlist
\item
  设 E、\(E_{1}\) 是两个 A-模，\(\Phi\) （相应地 \(\Phi_{1}\) ）是 E （相应地 \(E_{1}\) ）上的一个半双线性型。假设 \(\Phi_{1}\) 非退化，且存在一个元素 \(\alpha\in A\) 与 E 到 \(E_{1}\) 上的一个双射 u，使得 \(\Phi_{1}(u(x), u(y))=\Phi(x,y)\alpha\) （对 E 中一切的 x, y）。证明：\(1^{\circ}\Phi\) 非退化；\(2^{\circ}u\) 是线性的；\(3^{\circ}\) 若 \(E_{1}\) 是忠实 A-模，则 E 也是，且 \(\alpha\) 不是 A 中的右零因子；\(4^{\circ}\) 若 \(\Phi_{1}\) 在 A 中取的值都不是左零因子，则 \(\Phi\) 亦然。
\end{enumerate}

¶ 8) 设 A 是一个体，E₁、E₂ 是 A 上两个非零向量空间，Φ₁ （相应地 Φ₂）是 E₁ （相应地 E₂）上的一个非退化半双线性型（关于 A 的一个反自同构 J₁ （相应地 J₂））。设 u 是 E₁ 到 E₂ 上的一个线性映射，使得关系 Φ₁(x, y) = 0 蕴涵 Φ₂(u(x), u(y)) = 0。

\begin{enumerate}
\def\labelenumi{\alph{enumi})}
\item
  证明 \(u\) 是 \(\mathbf{E}_{1}\) 到 \(\mathbf{E}_{2}\) 上的一个双射。（若 \(u(0)\) 不化为 0，证明 \(\mathbf{E}_{1}\) 中会存在两个向量 \(a\) 、\(b\) ，使得 \(u(a) \neq 0\) 、\(u(b) = 0\) 且 \(\Phi_{1}(a, b) \neq 0\) ；若 H 是由 \(x \in \mathbf{E}_{1}\) （使 \(\Phi_{1}(a, x) = 0\) 者）所成的超平面，注意将有 \(u(\mathrm{H}) = \mathrm{E}_{2}\) 。）
\item
  证明：若 dim \(E_{1} \geqslant 2\) ，则存在 \(\alpha \in A\) ，使得有 \(\Phi_{2}(u(x), u(y)) = \Phi_{1}(x, y) \alpha\) （对 \(E_{1}\) 中一切的 x, y）。（对每个 \(y \in E_{1}\) ，证明存在一个元素 \(m(y) \in A\) ，使得 \(\Phi_{2}(u(x), u(y)) = \Phi_{1}(x, y)m(y)\) （对一切 \(x \in E_{1}\) ），且若 y 与 \(y'\) 在 \(E_{1}\) 中线性无关，则有 \(m(y + y') = m(y) = m(y')\) 。）
\end{enumerate}

\begin{enumerate}
\def\labelenumi{\arabic{enumi})}
\setcounter{enumi}{8}
\tightlist
\item
  设 A 是一个体，E、F 是 A 上两个左向量空间，\(\Phi\) 是 \(E \times F\) 上的一个非退化半双线性型（关于 A 的一个反自同构 J）。
\end{enumerate}

\begin{enumerate}
\def\labelenumi{\alph{enumi})}
\item
  设 M 是 E 的一个子空间，N 是 F 的一个子空间，使得 \(N \supset M^{0}\) 且 \(M \supset N^{0}\) 。证明：若空间 \(N/M^{0}\) 、\(M/N^{0}\) 之一是有限维的，则另一个也是，且这两个空间的维数相等。
\item
  设 M、\(M'\) 是 E 的两个子空间，使得 \(M^{00} = M\) 且 \(M'\) 是有限维的；证明我们有 \((M \cap M')^{0} = M^{0} + M'^{0}\) 与 \((M + M')^{00} = M + M'\) 。（把 a) 应用于子空间 \(M'\) 与 \(M^{0} + M'^{0}\) ，证明 \(\dim(M \cap M') = \text{codim}(M^{0} + M'^{0})\) ；把 a) 应用于子空间 \(M + M'\) 与 \(M^{0}\) ，证明
\end{enumerate}

\[
\dim ((M + M ^ {\prime}) ^ {0 0} / M) = \dim ((M + M ^ {\prime}) / M).
\]

\begin{enumerate}
\def\labelenumi{\alph{enumi})}
\setcounter{enumi}{2}
\item
  若 E = F 且 M 是 E 的一个子空间，使得 \(E = M^{0} + M^{00}\) ，证明 E 是 \(M^{0}\) 与 \(M^{00}\) 的直和。
\item
  设 E 是交换域 A 上的一个向量空间，具有可数无穷基 \((e_{n})_{n\geqslant0}\) ，并设 \(\Phi\) 是 E 上的对称双线性型，使得 \(\Phi(e_{n}, e_{n}) = 1\) （对一切 n），\(\Phi(e_{i}, e_{j}) = 0\) （对 \(i \geqslant 1\) 、\(j \geqslant 1\) 且 \(i \neq j\) ），且 \(\Phi(e_{0}, e_{n}) = 1\) （对一切 \(n \geqslant 1\) ）。证明 \(\Phi\) 非退化。设 M （相应地 N）是 E 的由诸 \(e_{2k}\) （相应地诸 \(e_{2k-1}\) ）（\(k \geqslant 1\) ）生成的子空间，并设 H = M + N，它是 E 的一个超平面。证明我们有 \(M^{0} = N\) 、\(N^{0} = M\) 、\(H^{00} = E \neq H\) 、\((M \cap N)^{0} \neq M^{0} + N^{0}\) 与 \((M + N)^{00} \neq M + N\) ，尽管 \(M^{00} = M\) 、\(N^{00} = N\) ；若 L 是由 \(e_{0}\) 与 \(e_{1}\) 生成的 2 维子空间，则有 \((L \cap H)^{0} \neq L^{0} + H^{0}\) 。
\end{enumerate}

¶ 10) 设 E、E′ 分别是体 A、A′ 上两个左向量空间，维数 ≥3；设 ℱ(E) （相应地 ℱ(E′)）是 E （相应地 E′）的有限维子空间所构成的（关于包含关系的）格序集。

\begin{enumerate}
\def\labelenumi{\alph{enumi})}
\item
  设 \(p\) 是 \(\mathfrak{F}(\mathrm{E})\) 到 \(\mathfrak{F}(\mathrm{E}')\) 中的一个映射，使得对每个 \(\mathbf{M} \in \mathfrak{F}(\mathrm{E})\) 有 \(\dim p(\mathbf{M}) = \dim \mathbf{M}\) ，且对 \((\mathbf{M}, \mathbf{N})\) （\(\mathfrak{F}(\mathrm{E})\) 的元素的对）有 \(p(\mathbf{M} + \mathbf{N}) = p(\mathbf{M}) + p(\mathbf{N})\) 。证明 \(p\) 是单射；若 \(p\) 是双射，则存在一个双射半线性映射 \(u\) （从 \(\mathrm{E}\) 到 \(\mathrm{E}'\) 中），使得有 \(u(\mathbf{M}) = p(\mathbf{M})\) （对一切 \(\mathbf{M} \in \mathfrak{F}(\mathrm{E})\) ）（利用第 II 章，第 2 版，附录 III，习题 10）。
\item
  给出一个例子，其中 \(A' = A\) 是交换的，\(E' = E\) 是有限维的，且存在一个 \(\mathfrak{F}(E)\) 到其自身的映射 p，使得 \(\dim p(M) = \dim M\) 、\(p(M + N) = p(M) + p(N)\) 、\(p(M \cap N) = p(M) \cap p(N)\) ，但不存在 E 到其自身的单射半线性映射 u，使得 \(u(M) = p(M)\) （对 \(M \in \mathfrak{F}(E)\) ）。（考虑存在 A 的一个有限次数且同构于 A 的扩体 \(A''\) 的情形，例如 \(\mathbf{A} = \mathbf{F}_{p}(\mathbf{X})\) ，其中 \(p\) 是素数；这时 \(\mathrm{E}'' = \mathrm{A}'' \otimes_{\mathbf{A}} \mathrm{E}\) 是 \(\mathrm{A}''\) 上的向量空间，其维数与 \(\mathrm{E}\) 在 \(\mathrm{A}\) 上的维数相同；考虑映射 \(\mathrm{M} \to \mathrm{A}'' \otimes_{\mathbf{A}} \mathrm{M}\) 。
\item
  我们假设 \(A' = A\) 。设 \(\omega\) 是 \(\mathfrak{F}(E)\) 到 \(\mathfrak{F}'(E')\) （\(E'\) 的有限余维数子空间所成之集）上的一个双射，使得 \(\text{codim } \omega(M) = \dim M\) ，且 \(\omega(M + N) = \omega(M) \cap \omega(N)\) 。证明存在一个半双线性型 \(\Phi\) （在 \(E \times E'\) 上，右非退化），且使得 \(\omega(M) = M^0\) （对一切 \(M \in \mathfrak{F}(E)\) ）。（利用第 II 章，第 4 节，第 6 小节，定理 1）。
\end{enumerate}

¶ 11) a) 设 A 是一个左、右阿廷环（第 VIII 章，第 2 节，第 3 小节）。证明下列条件等价：1° 每个左理想（相应地右理想）≠ A 的右零化子（相应地左零化子）不化为 0；2° 每个单左 A-模（相应地右 A-模）的对偶不化为 0；3° 每个有限生成左 A-模（相应地右 A-模）的对偶不化为 0。称环 A 满足条件 (N \(_{s}\) ) （相应地 (N \(_{a}\) )），如果它对左 A-模（相应地右 A-模）满足这些条件。

\begin{enumerate}
\def\labelenumi{\alph{enumi})}
\setcounter{enumi}{1}
\item
  设 A 是一个满足条件 (N \(_{a}\) ) 的左、右阿廷环，E （相应地 F）是 A 上具有可数基的自由左 A-模（相应地右 A-模），Φ 是 E × F 上的一个双线性型，使得 s \(_{\Phi}\) 是单射。设 M 是 E 的一个自由子模；证明存在 F 的一个自由子模 N 与 M （相应地 N）的基 (e \(_{n}\) ) （相应地 (f \(_{n}\) )），使得 Φ(e \(_{i}\) , f \(_{j}\) ) = δ \(_{ij}\) ；此外，对 e \(_{1}\) 可取 M 的任意自由元素。（注意若 x 是 M 的自由元素，则 F 在映射 y → Φ(x, y) 下的像是整个环 A；然后用归纳法构造基 (e \(_{n}\) ) 与 (f \(_{n}\) )。）由此导出：若 M 的基 (e \(_{n}\) ) 是有限的，则 E 是 M 与 N \(^{0}\) 的直和，且有 M \(^{00}\) = M。
\item
  我们保持 \(b\) 的假设，并再设 A 满足条件 \((\mathrm{N}_{s})\) 且 \(d_{\Phi}\) 是单射。证明这时 E 中存在一个基 \((e_{n})\) 且 F 中存在一个基 \((f_{n})\) ，使得 \(\Phi(e_{i}, f_{j}) = \delta_{ij}\) 。（用 \(b\) ，以归纳法交替地确定 \(e_{n}\) 与 \(f_{n}\) 。）
\end{enumerate}

¶ 12) 设 E 是一个可数型实希尔伯特空间，\(\Phi(x, y)\) 是 E 中的标量积。证明：不存在 R 上的向量空间 E 的两个代数基 \((e_{\lambda})\) 、\((f_{\mu})\) 组成的组，使得对每一对指标有 \(\Phi(e_{\lambda}, f_{\mu}) = \delta_{\lambda y}\) 。（首先注意这些基的指标集将具有连续统的势（《拓扑向量空间》，第 II 章，第 3 节，习题 15）；然后考虑 E 的一个（可数）正交基 \((a_n)\) ，并注意由诸 \(a_n\) 生成的子空间包含于 \((e_{\lambda})\) 的一个可数子族所生成的子空间中。
% semantic-fragments: legacy-input-seam
\relax{}
